% Generated draft: requires mathematical and source review.
\documentclass[11pt]{article}
\usepackage[margin=25mm]{geometry}
\usepackage{fontspec}
\setmainfont{Latin Modern Roman}
\usepackage{amsmath,amssymb,mathtools,unicode-math}
\setmathfont{Latin Modern Math}
\usepackage{enumitem,needspace,xurl,hyperref}
\setlength{\parindent}{0pt}
\setlength{\parskip}{5pt}
\setlength{\emergencystretch}{3em}
\newcommand{\N}{\mathbb{N}}
\newcommand{\Z}{\mathbb{Z}}
\newcommand{\Q}{\mathbb{Q}}
\newcommand{\R}{\mathbb{R}}
\newcommand{\C}{\mathbb{C}}
\newcommand{\F}{\mathbb{F}}
\newcommand{\A}{\mathbb{A}}
\newcommand{\PP}{\mathbb{P}}
\newcommand{\E}{\mathbb{E}}
\newcommand{\eps}{\varepsilon}
\newcommand{\dd}{\,\mathrm{d}}
\begin{document}
\section*{UnsolvedMath ID: 30001872}
\textbf{Source problem code:} catalogue entry
\subsection{Definitions and mathematical statement}
A line bundle \(L\) on a projective variety \(X\) is very ample if its global sections define a closed immersion into projective space.

An Okounkov body requires auxiliary data. In the usual flag construction, write \(n=\dim X\) and choose an admissible flag
\[
X=Y_0\supset Y_1\supset\cdots\supset Y_n=\{p\},
\]
where each \(Y_i\) is irreducible of codimension \(i\) in \(X\), and the flag members are nonsingular at \(p\). In a setting where this construction is defined, successive orders of vanishing along the flag, with division and restriction at each step, assign a vector \(\nu_{Y_\bullet}(s)\in\Z_{\geq 0}^{n}\) to each nonzero section \(s\in H^0(X,L^{\otimes m})\). The associated Okounkov body is
\[
\Delta_{Y_\bullet}(L)
=
\overline{\operatorname{conv}\!\left(
\bigcup_{m\geq 1}
\left\{\frac{\nu_{Y_\bullet}(s)}{m}:
0\neq s\in H^0(X,L^{\otimes m})\right\}
\right)}\subseteq\R^n.
\]
This recalls the flag construction conditionally; it does not supply missing hypotheses for the source question.

\textbf{Source question, with unresolved qualifications.} For every projective variety \(X\) and every very ample line bundle \(L\) on \(X\), does there exist a ``reasonable'' map from \(X\) to its Okounkov body, schematically
\[
f\colon X\longrightarrow\Delta_{Y_\bullet}(L)?
\]
The supplied record does not define ``reasonable,'' specify the base field or the meaning of the underlying domain of \(f\), or specify the admissible flag or alternative valuation data. It also does not state whether the auxiliary data are fixed beforehand or may be chosen together with \(f\). Consequently, this is an explicitly qualified question rather than a fully determined mathematical proposition. The statement asks for a map \emph{to} the body, whereas the supplied literature assessment discusses a map \emph{onto} it; surjectivity is therefore an unresolved requirement.

The supplied dated assessment describes constructions under finite-generation hypotheses and asks for a general construction beyond those hypotheses. It does not identify the object required to be finitely generated or the exact properties of those constructions. No finite-generation hypothesis is imposed on the universal source question.
\subsection{Short English statement}
The question asks whether every projective variety equipped with a very ample line bundle admits a suitable map to an associated Okounkov body. The intended suitability conditions, base field, auxiliary valuation data, and possible surjectivity requirement are unspecified. The supplied status is partially solved, based on dated curation requiring independent review.
\subsection{Sources}
\begin{itemize}
\item Local UnsolvedMath registry, permanent ID \(30001872\), source problem code \texttt{catalogue entry}, title: \emph{Canonical Maps from Varieties to Okounkov Bodies}. The supplied record lists proposed year \(2011\) and identifies no individual proposer.
\item Supplied source citation: \emph{Mini-Workshop: New Developments in Newton-Okounkov Bodies (2012)}, Oberwolfach Reports, pp.~2327--2363, DOI \texttt{10.4171/owr/2011/41}. Original report link: \url{https://doi.org/10.4171/owr/2011/41}. The citation's parenthetical year differs from the proposed year and the year appearing in the DOI; this discrepancy remains unresolved.
\item Supplied literature source: \emph{Okounkov bodies and toric degenerations}, \url{https://arxiv.org/abs/1001.4566}. No authors or precise result locator are supplied in the record.
\item Literature assessment dated \texttt{2026-08-21}: the supplied curation labels the problem partially solved and reports a reasonable map onto the Okounkov body under finite-generation hypotheses, while describing the general question as extending beyond those hypotheses. This is supplied dated curation, not an independently verified literature assessment.
\item The registry labels the statement \texttt{corrected\_verified} and reports that its mathematical content was checked against the official Oberwolfach report, with notation, context, and question form normalized. This verification claim is supplied registry metadata and has not been independently checked here.
\end{itemize}
\end{document}
